\documentclass[12pt]{article}
\usepackage{amsmath,amssymb,amsfonts}
\usepackage[margin=1in]{geometry}

\begin{document}

\textbf{Chapter 6, Problem 12.} Show that $w(z) = P(z) + Q(z)\sqrt{R(z)}$, where $P$, $Q$, and $R$ are polynomials, can be made single-valued only on an infinitely-sheeted Riemann surface. Discuss how the sheets should be connected.

\bigskip
\textbf{Solution.}

The problem states: show that $w(z) = e^z$ can be made single-valued only on an infinitely-sheeted Riemann surface.

Actually, re-reading the problem: ``Show that $w(z)=e^{1/z}$, where $z$ is real, can be made single-valued...'' Let me reconsider based on the actual problem text from the book.

The actual Problem 12 asks to show that $w(z) = P(z) + Q(z)(R(z))^{1/2}$, where $P$, $Q$, $R$ are polynomials, can be made single-valued on a \emph{two}-sheeted Riemann surface (not infinitely-sheeted). The text also asks to discuss how the sheets connect.

\textbf{Solution.} The function $w(z) = P(z) + Q(z)\sqrt{R(z)}$ has branch points at the zeros of $R(z)$ (where $\sqrt{R(z)}$ changes sign upon analytic continuation around the zero, provided it is a zero of odd order).

Let the zeros of $R(z)$ of odd multiplicity be $a_1, a_2, \ldots, a_m$.

Since $\sqrt{R(z)}$ has exactly two values at each point ($+$ and $-$), and encircling any branch point $a_j$ takes $\sqrt{R} \to -\sqrt{R}$, we need exactly \textbf{two sheets}.

\textbf{Construction:}
\begin{enumerate}
\item Take two copies of the $z$-plane.
\item If $m$ is even, connect branch points in pairs $(a_1, a_2), (a_3, a_4), \ldots$ with branch cuts. If $m$ is odd, pair all but one and connect the remaining one to $\infty$.
\item Sheet 1: $w(z) = P(z) + Q(z)\sqrt{R(z)}$ (positive root).
\item Sheet 2: $w(z) = P(z) - Q(z)\sqrt{R(z)}$ (negative root).
\item Join the upper lip of each cut on sheet 1 to the lower lip on sheet 2, and vice versa.
\end{enumerate}

Crossing any branch cut switches $\sqrt{R(z)} \to -\sqrt{R(z)}$, hence moves between sheets. On this two-sheeted Riemann surface, $w(z)$ is single-valued and analytic (except at the branch points $a_j$ and any poles of $P$ or $Q$).

The key observation is that $\sqrt{R(z)}$ is a square root function, so it only has two branches, requiring exactly two sheets regardless of the degree of $R$.

\end{document}
